\begin{theorem}[liouville]
\label{liouville}
\lean{ArithmeticFunction.liouville}
% at /home/runner/work/trough/trough/library/PrimeNumberTheoremAnd/IwaniecKowalskiCh1.lean:1737.0-1748.48
\leanok
Liouville function: $\lambda(n) = (-1)^{\Omega(n)}$.
\end{theorem}
\begin{proof}
\leanok
The Liouville function $\lambda(n)$ is defined as $(-1)^{\Omega(n)}$, where $\Omega(n)$ is the total number of prime factors of $n$ counted with multiplicity. This means that for each prime factor of $n$, we contribute a factor of $-1$ to the product, and the overall sign of $\lambda(n)$ depends on whether the total number of prime factors is even or odd. Thus, we have $\lambda(n) = (-1)^{\Omega(n)}$ by definition.
\end{proof}
